\section{Conclusão}
\label{sec:conclusion}

\subsection{Sumário das Descobertas}

Este trabalho comparou três bancos de dados de séries temporais (Apache IoTDB, InfluxDB e
TimescaleDB) em nove padrões de carga de trabalho IoT em três escalas de volume de dados,
utilizando uma infraestrutura de benchmark totalmente automatizada e publicamente reprodutível,
construída sobre Docker Compose e iot-benchmark~\cite{iotbenchmark}.

Três descobertas principais emergem do estudo.

\textbf{Nenhum banco de dados domina em todas as cargas de trabalho.} O IoTDB atinge
23,6\,M pts/s de vazão de escrita sequencial na escala grande (59$\times$ mais que o
TimescaleDB), mas requer aproximadamente 15\,GB de RAM e não oferece suporte a junções
relacionais. A consulta de subamostragem \texttt{time\_bucket()} do TimescaleDB permanece
abaixo de 1\,ms em todo o crescimento de cem vezes no volume de dados, tornando-o a escolha mais
forte para análises em painel, enquanto seu caminho de escrita baseado em WAL limita a vazão
de carregamento em lote. O InfluxDB ocupa um meio-termo prático: melhor vazão de escrita do
que o TimescaleDB e melhor eficiência de memória do que o IoTDB, mas mais fraco em consultas
de alertas por limite e consultas de último valor em escala.

\textbf{As escolhas arquiteturais determinam o comportamento de escalonamento, não o ajuste.}
A vantagem de filtro de valor do IoTDB cresce com o tamanho do conjunto de dados
(crescimento de latência de 5,2$\times$ para dados cem vezes maiores) porque suas
estatísticas de mínimo/máximo por bloco fornecem melhor eliminação no nível de bloco à medida
que a densidade de dados aumenta. A latência de último valor do InfluxDB cresce linearmente
com os arquivos armazenados porque o motor TSM não tem atalho de último valor em memória.
A subamostragem do TimescaleDB é genuinamente de latência constante por escalas porque as
consultas \texttt{time\_bucket()} varrem uma janela de tempo fixa, não o conjunto de dados
completo. Esses comportamentos derivam de escolhas arquiteturais centrais de cada sistema e são,
portanto, mais estáveis do que ganhos obtidos por ajuste de configuração. Mudanças de
versão que alterem essas escolhas centrais, como o novo motor colunar do InfluxDB~3.0,
podem modificar os compromissos observados.

\textbf{A metodologia afeta criticamente a validade do benchmark.}
Duas conclusões amplamente declaradas das execuções iniciais foram revertidas por reexecuções
isoladas e ajustadas: a penalidade de 23\% fora de ordem do TimescaleDB na escala grande era
um artefato de despejo do cache de páginas (reexecução: 2,7\%), e sua saturação de CPU de
97,7\% reportada foi causada por competição de CPU compartilhada (reexecução: 77,4\%). As
medições de memória mudaram em uma ordem de magnitude quando o ciclo de vida compartilhado
de contêineres foi corrigido. Essas reversões demonstram que resultados de benchmark em
ambientes com recursos compartilhados requerem validação explícita por meio de reexecuções
isoladas antes de poderem ser apresentados como propriedades arquiteturais.

\subsection{Comparação Final}

A Tabela~\ref{tab:final} fornece uma comparação multidimensional dos três bancos de dados
nas dimensões mais relevantes para decisões de implantação IoT. As avaliações usam uma escala
de três níveis: \textbf{A} (forte), \textbf{M} (moderado), \textbf{F} (fraco).

\begin{table*}[t]
  \centering
  \caption{Comparação final do Apache IoTDB, InfluxDB e TimescaleDB nas dimensões mais
    relevantes para implantações IoT. \textbf{A}~=~forte, \textbf{M}~=~moderado,
  \textbf{F}~=~fraco. Valores de latência na escala pequena, configuração ajustada.}
  \label{tab:final}
  \begin{tabular}{lcccl}
    \hline
    \textbf{Dimensão} & \textbf{IoTDB} & \textbf{InfluxDB} & \textbf{TimescaleDB} & \textbf{Observações} \\
    \hline
    Vazão de escrita sequencial      & \textbf{A} & M & F & 2,2\,M / 397\,K / 174\,K pts/s (pequena) \\
    Vazão de carregamento em lote (L=1000) & \textbf{A} & A & M & 11,8\,M / 2,7\,M / 214\,K pts/s \\
    Escrita de ponto único (L=1)     & \textbf{A} & M & F & 28\,K / 5\,K / 5\,K pts/s \\
    Vazão fora de ordem (abs.)       & \textbf{A} & F & F & 1,87\,M / 255\,K / 160\,K pts/s; IoTDB lidera em cada escala \\
    Tolerância fora de ordem (rel.)  & M & F & \textbf{A} & $-$13\% / $-$36\% / $-$8\% de penalidade \\
    Latência de último valor         & A & \textbf{F} & \textbf{A} & 1,24 / 9,24 / 0,38\,ms; InfluxDB escala mal \\
    Latência de subamostragem        & A & M & \textbf{A} & 1,81 / 5,71 / 0,61\,ms; TimescaleDB constante \\
    Exportação de intervalo de tempo & A & M & \textbf{A} & 1,19 / 5,41 / 0,54\,ms \\
    Latência de filtro de valor      & \textbf{A} & M & F & 3,72 / 71,09 / 23,98\,ms; IoTDB sublinear \\
    Filtro de valor na escala grande & \textbf{A} & M & F & 19,4 / 496,8 / 510,0\,ms \\
    Footprint de memória (escrita)   & F & \textbf{A} & A & $\sim$1,7\,GB / $\sim$190\,MB / $\sim$239\,MB inicial \\
    Footprint de memória (estável)   & F & M & A & $\sim$10\,GB / 900\,MB / 1,8\,GB \\
    Suporte a SQL/relacional         & F & F & \textbf{A} & Dialeto SQL IoTDB; apenas Flux; PostgreSQL completo \\
    Ecossistema e ferramental        & M & M & \textbf{A} & Em crescimento; TS maduro; ecossistema PG completo \\
    Garantia de durabilidade (escrita) & M & M & \textbf{A} & Flush assíncrono; flush assíncrono; WAL síncrono \\
    Modelo de dados IoT (nativo)     & \textbf{A} & M & M & Hierarquia dispositivo/sensor; tags+campos; relacional \\
    \hline
  \end{tabular}
\end{table*}

A tabela suporta um framework de seleção baseado em carga de trabalho. Para pipelines IoT
dominados por escrita com RAM suficiente, o IoTDB é a escolha clara. Para cargas de trabalho
pesadas em análise com consultas de painel em janelas de tempo recentes, a compatibilidade SQL
do TimescaleDB e a latência constante do \texttt{time\_bucket()} tornam-no a opção mais
forte. Para cargas de trabalho de alertas por limite que devem consultar conjuntos de dados
históricos crescentes, o escalonamento sublinear de filtro de valor do IoTDB oferece uma
vantagem estrutural de longo prazo. Para cargas de trabalho mistas com volumes de escrita
moderados e memória restrita, o InfluxDB fornece o perfil mais equilibrado.

\subsection{Limitações}
\label{sec:limitations}

As seguintes limitações delimitam as conclusões deste estudo:

\textbf{Implantação em nó único.} Todos os bancos de dados foram executados em um único host,
eliminando a latência real de rede. O protocolo de linha HTTP do InfluxDB incorre em sobrecarga
significativa de ida e volta em implantações distribuídas que a latência de loopback
próxima de zero não captura.

\textbf{Testes de leitura sequenciais.} As cargas de trabalho de leitura são executadas
sequencialmente no mesmo conjunto de dados dentro de uma sessão. Testes posteriores se
beneficiam de caches de páginas do SO mais quentes e estruturas de banco de dados em memória.
O isolamento com qualidade de produção exigiria flush dos caches e reinicialização do banco
de dados entre os testes de leitura.

\textbf{Durabilidade diferida do IoTDB.} A vazão de escrita do IoTDB reflete o reconhecimento
em memória antes da persistência em disco. Implantações com requisitos estritos de durabilidade
precisariam configurar o modo WAL síncrono, o que pode reduzir a vazão.

\textbf{Aquecimento da JVM.} O desempenho de leitura e de escrita de ponto único do IoTDB
melhora significativamente ao longo da duração de uma sessão à medida que a compilação JIT
otimiza os caminhos quentes. Benchmarks curtos subestimam a vazão em estado estacionário.
Os valores na escala grande são os mais representativos do comportamento do IoTDB em produção.

\subsection{Trabalhos Futuros}
\label{sec:future-work}

Diversas direções permanecem abertas para investigação posterior:

\begin{itemize}
  \item \textbf{Implantações distribuídas.} Avaliar clusters multi-nó do IoTDB, InfluxDB
    Clustered e TimescaleDB com a extensão Citus quantificaria como o escalonamento horizontal
    afeta os compromissos escrita-versus-leitura observados aqui e se a sobrecarga de rede muda
    os custos relativos de protocolo.

  \item \textbf{Conjuntos de dados do mundo real.} Testar com conjuntos de dados industriais
    em larga escala (telemetria de veículos, dados de sensores de redes elétricas ou
    monitoramento de processos de manufatura em escala de bilhões de pontos) validaria se o
    escalonamento sublinear de filtro de valor do IoTDB continua se mantendo e revelaria novos
    padrões não visíveis nas cargas de trabalho sintéticas usadas aqui.

  \item \textbf{Overhead de durabilidade e alta disponibilidade.} Medir o custo de vazão e
    latência de replicação, failover e cenários de recuperação de falhas seria essencial para
    qualquer decisão de implantação em produção. O caminho de escrita assíncrono do IoTDB torna
    isso particularmente importante.

  \item \textbf{Evolução de versões.} Os três bancos de dados lançam atualizações com
    frequência. O IoTDB~2.x, InfluxDB~3.0 (com seu novo motor de armazenamento colunar baseado
    em Apache Arrow e DataFusion) e TimescaleDB~2.x introduzem mudanças arquiteturais que podem
    alterar materialmente os compromissos descritos neste estudo.
\end{itemize}

Os scripts de benchmark, as configurações Docker e os arquivos de dados produzidos por este
estudo estão publicamente disponíveis para facilitar a replicação e extensão por pesquisadores
futuros.
